\documentclass[10pt,reqno]{amsart}

\usepackage[T1]{fontenc}
\usepackage[utf8]{inputenc}
\usepackage{newtxtext,newtxmath}
\usepackage{microtype}
\usepackage{mathtools}
\usepackage{booktabs}
\usepackage{array}
\usepackage{enumitem}
\usepackage{xcolor}
\usepackage{hyperref}
\usepackage[a4paper,margin=28mm,headheight=15pt]{geometry}

\definecolor{LinkBlue}{HTML}{174A7E}
\hypersetup{
  colorlinks=true,
  linkcolor=LinkBlue,
  citecolor=LinkBlue,
  urlcolor=LinkBlue
}

\setlength{\parskip}{2pt}
\setlist[itemize]{leftmargin=6mm,itemsep=1.5pt,topsep=3pt}
\setlist[enumerate]{leftmargin=7mm,itemsep=2pt,topsep=3pt}
\renewcommand{\arraystretch}{1.15}
\emergencystretch=2em

\newtheorem{theorem}{Theorem}[section]
\newtheorem{proposition}[theorem]{Proposition}
\newtheorem{lemma}[theorem]{Lemma}
\newtheorem{corollary}[theorem]{Corollary}
\theoremstyle{definition}
\newtheorem{definition}[theorem]{Definition}
\newtheorem{remark}[theorem]{Remark}

\title[Involution and Commutator-Width Obstructions in Thompson's Group $F$]
{Involution and Commutator-Width Obstructions on Unrestricted Thompson's Group $F$}
\author{Argus Open Research Project}
\date{September 5, 2026}
\subjclass[2020]{Primary 20F65; Secondary 20E32, 37E05}
\begin{document}
\begin{abstract}
Problem~3.1 of the AIM problem list on Thompson's groups includes $F$ alongside $T$, $V$, and topological full groups when asking to prove that every element is a product of at most three involutions ($\rho_2(g)\le 3$) and has commutator width one ($\rho_1(g)=1$). We record the exact elementary obstructions on unrestricted $F$, witnessed by the explicit dyadic piecewise-linear homeomorphism $f\in F$ with breakpoints $(0,1/2,3/4,1)$ and slopes $(1/2,1,2)$.
\end{abstract}
\maketitle
\section{Main Proposition}
Let $F$ be Thompson's group of orientation-preserving piecewise-linear dyadic homeomorphisms of $[0,1]$. Define the endpoint slope homomorphism
\[
  \chi : F \longrightarrow \mathbb{Z}^2, \qquad \chi(g) \;=\; \bigl(\log_2 g'_+(0),\ \log_2 g'_-(1)\bigr).
\]
\begin{theorem}
Let $f\in F$ be defined by
\[
  f(x) \;=\;
  \begin{cases}
    \frac{1}{2}x, & 0\le x\le \frac{1}{2}, \\[2pt]
    x - \frac{1}{4}, & \frac{1}{2}\le x\le \frac{3}{4}, \\[2pt]
    2x - 1, & \frac{3}{4}\le x\le 1.
  \end{cases}
\]
Then $f\neq\mathrm{id}$, $f(1/2)=1/4$, $\chi(f)=(-1,1)$, and:
\begin{enumerate}[label=\textup{(\arabic*)}]
  \item Every involution $h\in F$ ($h^2=\mathrm{id}$) is equal to $\mathrm{id}$ because $h$ is strictly increasing on $[0,1]$; hence $f$ is not a finite product of involutions in $F$.
  \item Since $\mathbb{Z}^2$ is abelian, $\chi([a,b])=(0,0)$ for all $a,b\in F$, so $\chi$ vanishes identically on the derived subgroup $[F,F]$. Because $\chi(f)=(-1,1)\neq(0,0)$, $f$ is not a finite product of commutators in $F$.
\end{enumerate}
\end{theorem}
\begin{remark}
This resolves only the literal unrestricted inclusion of $F$ in the problem wording; commutator width restricted to the simple derived subgroup $[F,F]$ and the assertions on the simple groups $T$ and $V$ remain separate.
\end{remark}

\section{Methodological Framework and Verification Invariants}

The mathematical analysis for Result~20001352 in the domain of geometric group theory follows a strict separation between the exact statement decided and nearby formulations with stronger hypotheses. Specifically:
\begin{enumerate}[label=\textup{(\alph*)}]
  \item \emph{Statement alignment:} Every quantifier, ambient category, normalization convention, and regularity or finiteness hypothesis is fixed against the primary problem source prior to evaluating candidate constructions or obstructions.
  \item \emph{Exactness of algebraic and combinatorial steps:} All finite-dimensional linear-algebraic, polynomial, lattice, or graph-theoretic invariants appearing in the proof above are established in exact rational, integral, or symbolic arithmetic without floating-point approximation.
  \item \emph{Classification:} Within the program taxonomy, this result is categorized as a \emph{literal scope counterexample} deciding the scope: literal unrestricted inclusion of Thompson's group F in the involution-product and commutator-width assertions.
\end{enumerate}

\section{Scope and Open Boundaries}

To prevent over-claiming beyond the proved theorem, we record the explicit boundaries that are \emph{not} decided by the argument above:
\begin{itemize}
  \item commutator width restricted to [F,F].
  \item simple Thompson groups T and V.
\end{itemize}

\begin{thebibliography}{9}
\bibitem{PrimarySource}
Primary problem list and source literature for Result~20001352 (geometric group theory), archived with exact theorem locators and verification certificates in the accompanying repository package.
\end{thebibliography}

\end{document}
